\section*{Exercises}
\addcontentsline{toc}{section}{Exercises}
\subsection*{A. Theory}
\begin{exercise}\exerciseid{18.A1} Explain why convolution with a fixed kernel is a linear map. \end{exercise}
\subsection*{B. By Hand}
\begin{exercise}\exerciseid{18.B1} Build the $4\times 4$ circulant matrix corresponding to the kernel $(1,0,-1,0)$. \end{exercise}
\subsection*{C. Python / NumPy}
\begin{exercise}\exerciseid{18.C1} Implement circular convolution directly for vectors of length $N$ and compare with multiplication by a circulant matrix. \end{exercise}
\subsection*{D. Challenge}
\begin{exercise}\exerciseid{18.D1} Explain why diagonalizing convolution is useful for fast computation. \end{exercise}
